\chapter*{Agradecimientos}

En primer lugar, quiero agradecer especialmente a mis padres, Laura y Osvaldo, las personas por quienes elegí este camino y quienes sembraron en mí, desde muy chico, la semilla de la curiosidad por resolver problemas, algo que años más tarde entendí que es la motivación de esta carrera.

A mi hermano, Mati, con quien siempre pude contar cuando lo necesité y que me acompañó a lo largo de este camino (y más).

A mi abuela Celia, a quien le atribuyo principalmente mi motivación para seguir una carrera universitaria, ya que me mostró desde muy chico la importancia de la educación pública. A la memoria de Julio e Irma, quienes no llegaron a verme en esta etapa, pero supieron acompañarme en otros momentos y sin quienes no hubiera llegado hasta acá.

A mis amigos del primario y a los que se sumaron en el secundario, que me enseñaron cosas que no se aprenden en otras instituciones y a quienes les agradezco haberme acompañado a lo largo de esta carrera.

A Cindy, Juli, Alan, Agus y Mai, con quienes formamos un grupo increíble y supimos acompañarnos mutuamente a lo largo de cursadas, trabajos prácticos y exámenes. A todos mis compañeros con los que tuve la suerte de compartir mates a lo largo de la carrera. A esta hermosa Facultad, que siempre tiene las puertas abiertas, y a la comunidad del Departamento de Computación, de la cual me sentí parte desde el primer día.

A mi director, Pablo, quien me propuso este tema de tesis y me acompañó durante un año con mucha dedicación y pasión por este trabajo.

A la Universidad de Buenos Aires y a la educación pública y gratuita.

Gracias.